% =====================================================================
% Subject taxonomy: body paragraph for the Dataset section, plus the
% figure caption. Numbers recomputed by make_category_figures.py and
% verified against the released files.
% =====================================================================

\subsection{Subject Coverage}
\label{sec:taxonomy}

We retain the curriculum subjects that the examination bodies themselves publish rather than
introducing a taxonomy of our own. A category scheme invented for a paper invites the objection
that its classes are neither exhaustive nor mutually exclusive, and the empirical difficulty
axis of Section~\ref{sec:difficulty} already carries the information such a scheme would be
proposing to capture. The published subjects are, by contrast, the units against which
candidates study and against which the examinations are blueprinted, so a result reported by
subject is interpretable without further argument.

$5{,}088$ of the $10{,}198$ items carry a subject label. Coverage is not random. The source
metadata tags mock and practice material with a paper identifier such as \texttt{2026 Practice
Exam A} and tags recalled examination items with a subject, so the labelled items are exactly
the held-out half and the public half carries no subject at all. Any subject-level statement in
this paper therefore describes the held-out partition, and we mark it as such wherever it
appears. Two spelling variants in the source were merged, \texttt{Quantitative} into
\texttt{Quantitative Methods} and \texttt{Ethic} into \texttt{Ethical and Professional
Standards}, which restores CFA Level~I to the ten subjects the curriculum defines.

The labelled items span $26$ distinct subjects across two examinations and five stages, and no
subject is shared between CFA and FRM. Figure~\ref{fig:taxonomy} lays out the full structure.
Three features of it matter for interpreting results by subject. CFA Levels~I and~II examine the
same ten subjects at increasing depth, which gives the benchmark a controlled
same-subject-deeper-treatment contrast that most examination corpora cannot offer. CFA Level~III
replaces most of those subjects with the specialisation pathways introduced in the current
curriculum, so it is a different territory rather than a harder version of the same one. And
Ethical and Professional Standards is the only subject present at all three CFA levels,
supplying $417$ items in total, which makes it the natural place to look for whether a system's
professional judgment improves with the depth the curriculum expects.

FRM is organised differently. Its two parts share no subject, Part~I covering foundations,
quantitative analysis, markets and valuation, and Part~II covering the four risk classes plus
the investment-management application. The nine FRM subjects are the complete set the
examination defines, so unlike CFA the FRM subject coverage in \textsc{FinExam-10K} is exhaustive
rather than a sample.

\begin{figure*}[t]
\centering
\includegraphics[width=\textwidth]{figures/fig_taxonomy.pdf}
\caption{Subject coverage of \textsc{FinExam-10K}. The five panels are the examination stages and
the nodes within each are the curriculum subjects the examination body publishes for that stage,
with the subject count and labelled item count given in each header. Thirty-seven nodes cover
$26$ distinct subjects, because Levels~I and~II examine the same ten and Ethical and
Professional Standards recurs at Level~III. Subjects are taken from the source metadata rather
than assigned by us, and two spelling variants were merged. The $5{,}088$ labelled items are
exactly the held-out half, since the public mock and practice half is tagged with paper
identifiers rather than subjects, so this figure describes the composition of the held-out
partition and not of the benchmark as a whole.}
\label{fig:taxonomy}
\end{figure*}
